% Copyright (c) 2025 by Daniel R. Driver.
% !TEX encoding = UTF-8 Unicode
% !TEX TS-program = XeLaTeX

\documentclass[titlepage]{article}

% This document presumes a file structure and set of inputs that are
% available at: git@github.com:danieldriver/syllabi.git

\newcommand\policy{../policy}
\newcommand\incl{../includes}
\ProvidesFile{variables.tex}[2022/09/11 v2.3 -- Syllabus variables]

\usepackage{xspace} % make manual spaces (like \mycmd\ ) unnecessary
\usepackage{xifthen} % provides \isempty test

% variables for internal use
\newcommand\prof{}
\newcommand\pdegree{}
\newcommand\pphone{}
\newcommand\pemail{}
\newcommand\poffice{}
\newcommand\phours{}
%
\newcommand\ccode{}
\newcommand\ctitle{}
\newcommand\cseries{}
\newcommand\cversion{}
\newcommand\csemester{}
\newcommand\cmeetson{}
\newcommand\cmeetsat{}
\newcommand\cmeetsin{}
\newcommand\cwebsite{}
\newcommand\cdescrip{}
\newcommand\cprereqs{}
\newcommand\edobject{}

% in case of fully online, BTh, or undergraduate courses - https://tex.stackexchange.com/a/5896
\newif\ifonline
\newcommand\Int[2]{\ifonline#1\else#2\fi}
\newif\ifbth
\newcommand\bth[2]{\ifbth#1\else#2\fi}
\newif\ifUG
\newcommand\UG[2]{\ifUG#1\else#2\fi}

% commands for setting variables in the preamble
\newcommand\professor[2][PhD]{
  \renewcommand\pdegree{#1\xspace}
  \renewcommand\prof{#2\xspace}}
\newcommand\phone[1]{
  \renewcommand\pphone{\addfontfeatures{Numbers=Monospaced}#1\xspace}}
\newcommand\email[1]{
  \renewcommand\pemail{\href{mailto:#1}{#1}\xspace}}
\newcommand\officehours[2][AST Library/Teams]{
  \renewcommand\poffice{#1\xspace}
  \renewcommand\phours{\href{https://outlook.office.com/bookwithme/user/b97b147a2e5c43b4a0e200c2cb630b77@astheology.ns.ca?anonymous}{#2}\xspace}%
  }
%
\newcommand\coursecode[2][1.0]{
  \renewcommand\cversion{#1\xspace}
  \renewcommand\ccode{#2\xspace}}
\newcommand\coursetitle[2][]{
  \ifthenelse{\isempty{#1}}%
    {}% do nothing if #1 is empty, else:
    {\renewcommand\cseries{#1\\[1ex]}}
  \renewcommand\ctitle{#2\xspace}}
\newcommand\semester[1]{
  \renewcommand\csemester{#1\xspace}}
\newcommand\meets[3]{
  \newcommand\AM{\textsc{am}}
  \newcommand\PM{\textsc{pm}}
  \renewcommand\cmeetson{#1\xspace}
  \renewcommand\cmeetsat{\Int{From 9:00 \AM}{#2}\xspace}
  \renewcommand\cmeetsin{\Int{\href{https://smu.brightspace.com/d2l/login}{Brightspace}}{#3}\xspace}}
\newcommand\website[2][]{
  \renewcommand\cwebsite{
  \ifthenelse{\isempty{#1}}%
    {\href{https://#2}}% copy url if #1 is empty, else
    {\href{#1}}% use specified url and print #2
  {#2}\xspace}}
\newcommand\cdescription[2][BF 1001; RM 1000 or GTRS 6000]{
  \renewcommand\cprereqs{#1}
  \renewcommand\cdescrip{#2\par}}
\newcommand\objectives[1]{
  \renewcommand\edobject{#1\par}}


%\onlinetrue % \Int{true}{false}
%\UGtrue % \UG{true}{false} terms for undergraduate students
\coursecode[1.5\UG{-UG}{}]{HB 3115\UG{-UG}{}}% version is 1.0 unless \cmd[1.1]{}
\coursetitle{Exodus and the Life of Moses}% \cmd[optional series]{title}

% Taught in:
%   - Summer 2018
%   - Fall 2025

\professor{Daniel R. Driver}% terminal degree is PhD unless \cmd[DPhil]{}
\phone{902-425-7051}
\email{daniel.driver@astheology.ns.ca}
\officehours{Thursdays, 1:00--3:00 \PM}

\semester{Fall 2025}
\meets{Mondays}% \meets{on}{at}{in} - on could be: \Int{Mondays}{Thursdays}
      {10:00 \AM--12:30 \PM\ Atlantic}% if online, it prints "From 9:00 AM"
      {MS Teams}% if online, it prints "Brightspace" with link
\website{danieldriver.com/courses}
\cdescription{% \cmd[standard prereqs]{from the current Academic Calendar}

Who is Moses? How does his story shape Israel’s story in Egypt and at
Sinai? This course will examine the life of Moses in the context of the
book of Exodus, with special attention to the book’s canonical shape and
historic reception. Students will also become familiar with some of the
current debate about the formation of the Pentateuch, particularly as it
relates to different accounts of Israel’s origins in the Hebrew
Bible/Old Testament.

}% end of course description
\objectives{%

By the end of the course students should be able to demonstrate a
critical awareness of the main contours of the Book of Exodus and, to a
more limited extent, of its modern interpretation and ancient reception
history. They should be able to communicate their own understanding of
who Moses is, with an ecumenical outlook that includes some knowledge of
Jewish and Christian tradition. They should develop a critical
self-awareness of their own patterns of thought about biblical story and
biblical law, particularly where these elements converge on the figure
of Moses. Ideally, they will discover in the study of Old Testament
scripture a genuine love of learning, all due humility before the task
of interpretation, and the courage to address themselves to text and
tradition.

Students should also be able to:
	summarize and understand major hermeneutical approaches to the Pentateuch;
	identify, analyze, and explicate short passages from it, especially from Exodus;
	appreciate, compare, and evaluate modern and traditional religious understandings of the text;
	generate and defend an original, coherent thesis about Moses.

% A Model of Learning Objectives
% Adapted from http://www.celt.iastate.edu/teaching/RevisedBlooms1.html
%
% A learning objective should contain a *subject* (a verb related to
% cognitive process) and an *object* (a noun identifying the knowledge
% or skill sought). Statements should be precise and measurable, e.g.,
% "By the end of this course, students will be able to <verb> <noun>."
%
% Anderson and Krathwohl name four *classes of knowledge* that students
% can be expected to acquire. Terms move from the concrete to the
% abstract:
%
%   1. factual           (basic elements, details, terminology)
%   2. conceptual        (interrelationship of basic elements)
%   3. procedural        (skills, methods, techniques, criteria)
%   4. metacognitive     (strategy, context and conditions, self-knowledge)
%
% Metacognitive knowledge is a special case. It concerns one's own
% "cognition and about oneself in relation to various subject matters"
% (Anderson and Krathwohl, 2001, p.44).
%
% A&K also identify 19 *cognitive processes*, arranged in six categories
% (a 2001 revision of Bloom's 1957 taxonomy) that progress from lower
% order to higher order thinking skills (from LOTS to HOTS):
%
%   1. remember
%      * recognizing     (identifying)
%      * recalling       (retrieving)
%   2. understand
%      * interpreting    (clarifying, paraphrasing, representing, translating)
%      * exemplifying    (illustrating, instantiating)
%      * classifying     (categorizing, subsuming)
%      * summarizing     (abstracting, generalizing)
%      * inferring       (concluding, extrapolating, interpolating, predicting)
%      * comparing       (contrasting, mapping, matching)
%      * explaining      (constructing models)
%   3. apply
%      * executing       (carrying out)
%      * implementing    (using)
%   4. analyze
%      * differentiating (discriminating, distinguishing, focusing, selecting)
%      * organizing      (finding coherence, integrating, outlining, parsing, structuring)
%      * attributing     (deconstructing)
%   5. evaluate
%      * checking        (coordinating, detecting, monitoring, testing)
%      * critiquing      (judging)
%   6. create
%      * generating      (hypothesizing)
%      * planning        (designing)
%      * producing       (constructing)
%
% If you reference a List of Measurable Verbs Used to Assess Learning
% Outcomes, remember that A&K reverse Bloom's 6th and 5th categories.

}% end of learning objectives

\ProvidesFile{preamble.tex}[2014/09/02 v1.1 -- Syllabus preamble]

% basic typography
\usepackage{fontspec,realscripts}
\defaultfontfeatures{Ligatures=TeX}
\setmainfont[%
  Numbers={Proportional,OldStyle},
  SmallCapsFeatures={RawFeature=+c2sc;+smcp,LetterSpace=1}]{Scala OT}
\newfontfamily\lining[Numbers={Lining}]{Scala OT}
\setsansfont[%
  SmallCapsFeatures={RawFeature=+c2sc;+smcp,LetterSpace=1}]{Scala Sans Pro}
\setmonofont[Scale=MatchLowercase]{Menlo}
\usepackage{sectsty}
\allsectionsfont{\sffamily\scshape}
\frenchspacing
\setlength{\emergencystretch}{3em} % prevent overfull lines

% custom font size and leading
\renewcommand\tiny{\fontsize{6}{9}\selectfont}
\renewcommand\scriptsize{\fontsize{7}{10}\selectfont}
\renewcommand\footnotesize{\fontsize{8}{10}\selectfont}
\renewcommand\small{\fontsize{9}{11}\selectfont}
\renewcommand\normalsize{\fontsize{10}{12}\selectfont}% base size
\renewcommand\large{\fontsize{11}{13}\selectfont}
\renewcommand\Large{\fontsize{12}{14}\selectfont}
\renewcommand\LARGE{\fontsize{16}{18}\selectfont}% "course syllabus \\ semester" benefits from more lead
\renewcommand\huge{\fontsize{20}{22}\selectfont}
\renewcommand\Huge{\fontsize{24}{26}\selectfont}

% layout packages: page, logo, tables
\usepackage[scale={0.6,0.8},
            xetex]{geometry}
\usepackage{graphicx}
\usepackage{array}     % allow insertions of column styling with >{}
\usepackage{booktabs}  % elegant horizontal rules in tables

% custom macros for a session count in the schedule of readings
\newcounter{session}
\newcounter{columns}
\newcounter{spancolumns}
\newcounter{midcolumns}
\newcounter{courseunit}
\newcommand\setcolumncount[2][0]{ % optionally set count to other than 0,
  \setcounter{session}{#1}        % e.g. to -1, or to a standing count
  \setcounter{columns}{#2}
  \setcounter{spancolumns}{#2-1}
  \setcounter{midcolumns}{#2-2}}
\newcommand\sessioncount{\stepcounter{session}\arabic{session}}
\newcommand\sessionskip[1]{\multicolumn{1}{@{}r@{ }}{#1}}
\newcommand\unit[1]{\multicolumn{\thecolumns}{c}{%
  \scshape\stepcounter{courseunit}\roman{courseunit}. #1}}
\newcommand\noclass[1]{\sessionskip{---.}&\multicolumn{\themidcolumns}{@{}l}{\itshape No Class: #1}}
\newcommand\rarr{\Large\char"EA31\hspace*{0.5pt}} % pointer finger w/ reminder tie in Scala OT
\newcommand\reminder[2]{\sessionskip{\rarr}&\multicolumn{\thespancolumns}{@{}l}{#1 \hfill #2}}

% color to match Tyndale's branding
\usepackage[usenames]{xcolor}
% predefined: black, white, red, green, blue, cyan, magenta, yellow
\definecolor{TyndaleURLs}{HTML}{0062A0} % links on tyndale.ca
\definecolor{TyndaleBlue}{cmyk}{1,1,0,.32}
\definecolor{TyndaleGold}{cmyk}{0,.27,1,0}
\definecolor{TyndaleRed}{cmyk}{0,1,.99,.04}
\definecolor{TyndaleBlack}{cmyk}{0,0,0,1}
\definecolor{TyndaleGreen}{cmyk}{.45,0,1,.24}
\definecolor{TyndaleOrange}{cmyk}{0,.79,1,0}
\definecolor{TyndaleAqua}{cmyk}{.47,0,.24,0}
\definecolor{TyndaleYellow}{cmyk}{.03,.03,.35,0}

% metadata (assumes a host of definitions are made in the main file)
\usepackage[setpagesize=false,     % leave this to geometry
            hyperfootnotes=false,  % fragile and distracting
            xetex]{hyperref}
\hypersetup{breaklinks=true,       % allow link text to break across lines
            colorlinks=true,       % colorlinks resets pdfborder to 0 0 0
            urlcolor=TyndaleURLs,  % for external links
            linkcolor=TyndaleRed,  % for normal internal links
            citecolor=TyndaleGold, % for bibliographical citations in text
            pdfauthor={\prof},
            pdftitle={\ccode: \ctitle},
            pdfsubject={Tyndale UC, \csemester},
            pdfcreator={github.com/danieldriver/syllabus}}
\urlstyle{same}                    % don't use monospace font for urls

% custom footlines
\usepackage{fancyhdr}
\pagestyle{fancy} % turn it on
\fancyhf{}        % reset everything
\renewcommand{\headrulewidth}{0pt} % remove header line as well
\lfoot{\sffamily\scshape\footnotesize\ctitle, \csemester}
\rfoot{\sffamily\scshape\footnotesize\prof\quad\thepage}

% language support: polyglossia needs to go last
\usepackage{polyglossia}
\setdefaultlanguage{english}
\setotherlanguage{hebrew}
\newfontfamily\hebrewfont[Script=Hebrew]{SBL BibLit}

% protect positioning of margin table in policy/grades -- needs to come after polyglossia
\usepackage{marginfix}

% gratuitous with custom title page, but useful as a fallback
\title{\ccode: \ctitle}
\author{\professor}
\date{\semester}


\begin{document}
\ProvidesFile{title.tex}[2013/09/06 v1.0 -- Syllabus title page]

\begin{titlepage}
  \begin{center}

    \begin{minipage}{\textwidth}
      \parbox[t]{0.5\textwidth}{
        \mbox{}\\[-\baselineskip] % dummy line to align parboxes
        \includegraphics[width=0.5\textwidth]{.syllabus/includes/TyndaleUC}}
      \hfill
      \parbox[t]{0.4\textwidth}{
        \raggedleft\LARGE\sffamily Course Syllabus\\
        \semester}
    \end{minipage}

    \vfill

    {\sffamily\textsc{\LARGE\MakeLowercase\coursecode}\\[2ex]
      \bfseries\itshape\Huge\coursetitle}

    \vfill

    \begin{tabular}{>{\bfseries}rl>{\bfseries}rl}
      \toprule
      Instructor & \professor                   & Course  & \liningnums{\coursecode}~\texttt{(v\version)} \\
      \midrule
      Phone      & \texttt{\phone}              & Meets   & \meeton                                       \\
      Email      & \href{mailto:\email}{\email} & Time    & \meetat                                       \\
      Office     & \officenumber                & Room    & \meetin                                       \\
      Hours      & \officehours                 & Website & \href{http://\website}{\website}              \\
      \bottomrule
    \end{tabular}

    \vfill

    \begin{description}
      \item[Commuter Hotline]
        Class cancellations due to inclement weather or illness will
        be announced on the commuter hotline at \texttt{416.226.6620
        x2187}. Alternately, weather cancellation information is posted
        at \href{http://tyndale.ca/weather}{tyndale.ca/weather}.
      \item[MyTyndale.ca]
        This course may have materials stored on its website, such as
        handouts or readings that may be needed in order to complete
        assignments. Students are responsible for checking these course
        pages on a regular basis. Here, too, students are able to view
        their grades throughout the semester. For more information see
        Section~\ref{mytyndale}, below.
      \item[Mail]
        Students are responsible for information communicated through
        their campus mailboxes and student e-mail accounts. A mailbox
        directory hangs beside the mailboxes. For more information
        contact the Registrar's office.
    \end{description}

  \end{center}

  \section{Course Description}
  \label{description}

  \emph{From the Academic Calendar:} \cdescription

\end{titlepage}
\setcounter{page}{2} % count the title page as page 1


\section{Learning Objectives}
\label{objectives}

\edobject

\section{Required Texts \& Materials}
\label{texts}

The following texts are required. Students should have their own copies.
Library copies that are not reference works have been placed on reserve.

\begingroup
\renewcommand{\section}[2]{}% temporarily remove the section heading
\begin{thebibliography}{Sommer}% use the longest item in the bibliography

	\bibitem[Nyssa]{nyssa} Gregory of Nyssa.
	\emph{The Life of Moses}. Translated by Abraham J. Malherbe and Everett Ferguson.
	New York: Paulist, 1978, rev. 2002; repr. New York: HarperOne, 2006.
	ISBN 978-0809121120 (Paulist) or 978-0060754648 (HarperOne).

	\bibitem[Sommer]{sommer} Benjamin D. Sommer.
	\emph{Revelation and Authority: Sinai in Jewish Scripture and Tradition}.
	New Haven: Yale University Press, 2015.
	ISBN 978-0300234688.

	\bibitem[von Rad]{vonrad} Gerhard von Rad.
	\emph{Moses}. 2nd ed. Edited by K.\,C. Hanson.
	Eugene, OR: Cascade Books, 2011.
	ISBN 978-1606087718.

\end{thebibliography}
\endgroup

Students should also have a good, modern translation of the Bible such
as the NRSV or NJPS. If you want a study Bible, I recommend either
Michael Coogan et al., eds., \emph{The New Oxford Annotated Bible: NRSV
with Apocrypha} (5th ed.; Oxford: OUP, 2018) or, with some superb notes
and essays but neither Apocrypha nor NT, Adele Berlin and Marc Zvi
Brettler, eds., \emph{The Jewish Study Bible: Second Edition} (2nd ed.;
Oxford: OUP, 2014).

\section{Supplementary Texts}
\label{supplementary}

Supplementary readings may be recommended throughout the semester.
Excerpts from this literature will either be placed on reserve or made
available through the course website. The following secondary titles are
notable but not required:

\begingroup
\renewcommand{\section}[2]{}% temporarily remove the section heading
\begin{thebibliography}{Childs}% use the longest item in the bibliography

	\bibitem[Buber]{buber} Martin Buber.
	\emph{Moses}.
	Oxford: East and West Library, 1946.
	
	\bibitem[Childs]{childs} Brevard S. Childs.
	\emph{The Book of Exodus: A Critical, Theological Commentary}.
	Louisville: Westminster John Knox Press, 1974.

	\bibitem[Schmid]{schmid} Konrad Schmid.
	\emph{Genesis and the Moses Story: Israel’s Dual Origins in the Hebrew Bible}. Translated by James D. Nogalski.
	Winona Lake, IN: Eisenbrauns, 2010.

\end{thebibliography}
\endgroup%

Also, the following basic works are worth consulting and even owning.
\cite{rlgs} in particular contains sound advice on core skills like
reading religious texts, writing essays and reviews, revising essays,
making oral presentations, and learning languages.

\begingroup
\renewcommand{\section}[2]{}% temporarily remove the section heading
\begin{thebibliography}{9}% use the longest item in the bibliography

	\bibitem[Making]{rlgs} Northey, Margot, Bradford A. Anderson, and Joel N. Lohr.
	\emph{Making Sense in Religious Studies: A Student's Guide to Research and Writing}.
	3rd ed. Toronto: Oxford University Press, 2019. AST Library: Ref BL 41.N67 2019.

	\bibitem[SBL2]{sbl2} Collins, Billie Jean, et al.
	\emph{The SBL Handbook of Style}.
	2nd ed. Atlanta: SBL Press, 2014.
	Designed to augment \href{https://proxy.openathens.net/login?qurl=https%3A%2F%2Fwww.chicagomanualofstyle.org%2Fbook%2Fed17%2Ffrontmatter%2Ftoc.html}{\emph{Chicago Style}}
	(the standard at AST), there is also a free
	\href{https://www.sbl-site.org/wp-content/uploads/2025/04/SBLHSsupp2015-02.pdf}{Student Supplement for SBL2}. AST Library: Ref PN 147.S26 2014.

\end{thebibliography}
\endgroup

\section{Course Outline}
\label{outline}

% For ease of citation:
\newcommand\bsc[1]{\cite[#1]{childs}}
\newcommand\bds[1]{\cite[#1]{sommer}}
\newcommand\gvr[1]{\cite[#1]{vonrad}}
\newcommand\nys[1]{\cite[#1]{nyssa}}

We will adhere to the schedule in \autoref{schedule} as closely as
possible, though the professor reserves the right to adjust it to suit
the needs of the class.

\setcolumncount{4}% set up \sessioncount, \unit{}, \noclass{}, and \reminder{memo}{date} macros
\begin{table}[htb]% set to `p' to put the schedule on its own page
  \centering
  \begin{tabular}{>{\sessioncount.}r@{ }llr}% make sure the column config agrees with \setcolumncount
	\toprule
	\sessionskip{\textbf{\S}.}&\textbf{Primary}&\textbf{Secondary}&\textbf{Due}\\
	\midrule

		& Exodus 1                         &               &  8 Sep. \\ % or \Int{14}{12} Sep.
		& Exodus 2                         & \nys{1–68}    & 15 Sep. \\
		& Exodus 3                         & \nys{69–132}  & 22 Sep. \\
		& Exodus 7                         & \gvr{1–44}    & 29 Sep. \\
		& Exodus 12                        & \gvr{45–87}   &  6 Oct. \\ [1ex]
	\reminder{The first paper is \textbf{due} before the \textbf{sixth week} of class.}{26 Oct.} \\
		
	\noclass{Thanksgiving}                                 & 13 Oct. \\
	\noclass{Reading Week}                                 & 20--24 Oct. \\ [1ex]

		& Exodus 16                        & \bds{1–26}    & 27 Oct. \\
		& Exodus 19                        & \bds{27–98}   &  3 Nov. \\
		& Exodus 20                        & \bds{99–146}  & 10 Nov. \\
		& Exodus 24                        & \bds{147–187} & 17 Nov. \\
		& Film: \emph{The Prince of Egypt} &               & 24 Nov. \\
		& Exodus 32--33                    & \bds{188–208} &  1 Dec. \\
		& Exodus 33--34                    & \bds{209–251} &  8 Dec. \\ [1ex]

	\reminder{End of Term: Final marks are due for this course.}{15 Dec.} \\

	\bottomrule
  \end{tabular}
  \caption{Schedule of Readings}
  \label{schedule}
\end{table}

See the AST website for a list of other \href{https://www.astheology.ns.ca/students/index.html}{important dates}.

\section{Evaluation}
\label{evaluation}

The grade structure for \ccode has the following elements.

\begin{enumerate}

	\item To facilitate regular engagement with the biblical text,
	\textbf{notes and quotes} will be solicited from students at the
	start of each class. Prepare by giving careful attention to the
	primary reading in a modern translation (such as NJPS or NRSV).
	Annotate it. Mark up your Bible. Make excursive notes as you read.
	Reflect: What do you note about the material? What words, phrases,
	or ideas stand out to you? \UG{Undergraduate students should augment
	their observations with the footnotes in a good Study
	Bible.}{Graduate students must augment this work by consulting a
	unique commentator, to be changed midway through the course to a
	commentator from a different era.} All students should come prepared
	to share and discuss your observations with the class.

	\item Imagine that a church press is releasing a new edition of a
	classic treatise on Moses, and you have agreed to write a new
	preface. Compose a new \textbf{book preface} for either \cite{nyssa}
	or \cite{vonrad}. Describe what your selected book contributes to a
	Christian understanding of Moses. Be sure to touch on aspects of
	biography and authority (legal and moral). Contextualize your
	description by comparing your book to a Christian portrait from a
	very different era, whether modern (20th C. German in \cite{vonrad})
	or ancient (4th C. Cappadocian Greek in \cite{nyssa}). Focus on one
	work and use the other for auxiliary support. By all means “spoil”
	the works, and illustrate your account with abundant paraphrase and
	a few well-chosen quotes. Your preface should be about
	\UG{2,000}{3,000} words long. It is due before the sixth week of
	class.

	\item \textbf{Teach a chapter} from \cite{sommer}. In most cases,
	students will work in pairs. Undergraduate students should pair up
	with graduate students. One or two graduate students may elect to
	work alone. Teaching sessions should last 30 minutes, with about 20
	min given to presentation and 10 min to class discussion.

	\item A \textbf{final thesis} gives students an opportunity to
	develop their own answer to the question “Who is Moses?” The paper
	must have a clear thesis, which should govern the entire discussion.
	The thesis needs to be supported by exegesis of passages form Exodus
	and research into the literature. The paper should also show an
	awareness of the readings, lectures, and classroom discussions. It
	should be \UG{2,500}{4,000} words long, the first 500 of which are
	to be read aloud at the start of the last class. Let this
	introduction feature a single-sentence thesis statement. The paper
	is due by the last class.

\end{enumerate}

For guidance on how to approach these and other assignments, see
\cite{rlgs} in \autoref{supplementary}, above. The breakdown for the
semester's total work is shown in \autoref{grade-dist}.

\begin{table}[htbp]
  \centering
  {\lining
  \begin{tabular}{lr}
    \toprule
    Notes \& Quotes  & 20\% \\
    Book Preface     & 25\% \\
    Teaching Session & 25\% \\
    Final Thesis     & 30\% \\
    \bottomrule
  \end{tabular}}
  \caption{Distribution of Grades}
  \label{grade-dist}
\end{table}

\ProvidesFile{grades.tex}[2023/09/06 v2.5 -- Course policy]

AST's \href{https://www.astheology.ns.ca/students/resources.html}{Academic
Calendar} provides guidelines and criteria for academic assessment. Marks are
assigned by letter grade using these benchmarks. Note that graduate and
undergraduate programs are currently assessed on a different scales at AST (see
\autoref{grade-syst}). For graduate students, a grade of C is deemed
``Unsatisfactory,'' whereas for undergraduate (BTh) students, a C is
``Satisfactory'' and the grade of D is possible for ``Marginal'' performance.
Other grades have different definitions, too.

\begin{table}[htbp]
 \centering
 {\lining\footnotesize
 \begin{tabular}{lll|ll}
   \toprule
               & \multicolumn{2}{c|}{Graduate Courses}   & \multicolumn{2}{c}{Undergraduate Courses} \\
               \cline{2-5}
   A+          & 94--100    & Exceptional               & 90--100    & Excellent                 \\
   A           & 87--93     & Outstanding               & 85--89     & Excellent                 \\
   A\char"2212 & 80--86     & Excellent                 & 80--84     & Excellent                 \\ [1ex]
   B+          & 77--79     & Good                      & 75--79     & Good                      \\
   B           & 73--76     & Acceptable                & 70--74     & Good                      \\
   B\char"2212 & 70--72     & Marginally Acceptable     &            &                           \\ [1ex]
   C           & 60--69     & Unsatisfactory            & 60--69     & Satisfactory              \\
   D           &            &                           & 50--59     & Marginal                  \\
   F           & 0--59      & Failure                   & 0--49      & Failure                   \\
   FP          & 0          & Failure due to Plagiarism & 0          & Failure due to Plagiarism \\
   \bottomrule
 \end{tabular}}
 \caption{Outline of Grades at AST}
 \label{grade-syst}
\end{table}

\ProvidesFile{other.tex}[2018/0/03 v2.4 -- Course policy]

\section{Other Course Policy}
\label{policy}

Late work will not be accepted, except in genuinely extenuating
circumstances. Students must submit something before the deadline if
they wish to receive credit. Unless I state otherwise, written
assignments are to be uploaded by 11:59 \PM\ on the date indicated.

Essay submissions must be typewritten, double-spaced, and formatted as
PDFs. They should be free from error. In this course they should follow
either \emph{The SBL Handbook of Style} (refer to the online
\href{https://www.sbl-site.org/assets/pdfs/pubs/SBLHSsupp2015-02.pdf}{
Student Supplement}), or \emph{The Chicago Manual of Style}, per AST's
policy for major papers. As a reminder, AST also upholds an Inclusive
Language Policy.

\href{http://www.eerdmans.com/Pages/Item/59043/Commentary-Statement.aspx}{Plagiarism},
if \href{https://www.theguardian.com/world/2013/feb/09/german-education-minister-quits-phd-plagiarism}{detected},
will result in failure of the course.

Students should request permission to record a class or lecture. If
permission is granted, or if recordings are provided (as in the case of
an online course), I stipulate that all recordings be for personal use
only. They may not be shared or distributed.

If you have abilities or disabilities that require modifications to the
assessment process or other aspects of this course, please advise the
course instructor as soon as possible.

Finally, I encourage the conscientious use of laptops, tablets, and
other technology in my classes. In classroom settings, realize that, as
\href{http://dx.doi.org/10.1016/j.compedu.2012.10.003}{cognitive
psychologists have demonstrated}, ``laptop multitasking hinders
classroom learning for both users and nearby peers.'' Do your part to
foster an environment of open dialogue by honouring the presence of your
classmates. In online settings, consider both the physical environment
in which you choose to work and the virtual environment that you help
create through your participation in various forums. Let your engagement
in this course be marked by rigour and charity alike.

\section{Further Bibliography}
\label{bibliography}

For additional literature, I recommend exploring \href{https://go.openathens.net/redirector/astheology.ns.ca?url=https://www.oxfordbibliographies.com/obo/page/biblical-studies}{Oxford Bibliographies: Biblical Studies (Full Text)}.
You can access the database automatically while on campus or remotely
with OpenAthens credentials. Among the many articles by subject-area
specialists, see especially these three by Thomas B. Dozeman: “Moses,”
“Pentateuch,” and “Book of Exodus.”

\end{document}
